% ============================================================
% CAPÍTULO 2 — Logros del Proyecto
% ============================================================
\chapter{Logros del Proyecto}
\label{cap:logros}

\section{Hitos t\'ecnicos alcanzados}

El sistema cubre las cuatro capas proyectadas y un SDK cient\'ifico
independiente. A continuaci\'on se enumeran los hitos verificados por
capa.

\subsection{Capa de Campo (Firmware ESP32)}
\begin{itemize}[leftmargin=*]
    \item Implementaci\'on completa de los tres roles (Sensor, Actuador,
    Gateway) con c\'odigo compartido y comportamiento configurable en
    compilaci\'on.
    \item Protocolo binario Demeter v2 funcional: \textbf{14} tipos de
    mensaje operativos, handshake TCP-like (SYN / SYN-ACK / ACK),
    gesti\'on de reintentos.
    \item Operaci\'on en \textit{Deep Sleep} para nodos sensor
    ($\sim$80\,mA en activo, $\sim$0.01\,mA en reposo).
    \item Tabla de rutas din\'amica con comando \texttt{ROUTE\_ADD}.
    \item Pipeline de validaci\'on por capas: SYNC + LENGTH + CRC + DST\_ID.
\end{itemize}

\subsection{Capa de Borde (Raspberry Pi)}
\begin{itemize}[leftmargin=*]
    \item Servicio Python \texttt{asyncio} con cuatro tareas concurrentes.
    \item Decodificaci\'on bidireccional Protocolo Demeter $\leftrightarrow$ JSON.
    \item Cliente WebSocket TLS persistente con reconexi\'on autom\'atica.
    \item Containerizaci\'on Docker con mapeo directo del UART f\'isico.
    \item Cach\'e local de topolog\'ia (\textit{Device Manager}) persistente.
\end{itemize}

\subsection{Capa de Servidor (Google Cloud)}
\begin{itemize}[leftmargin=*]
    \item Stack de seis contenedores en red Docker privada
    (\texttt{demeter-net}).
    \item Backend FastAPI con \textbf{31} endpoints REST (19 GET, 10 POST,
    1 PATCH, 1 DELETE) y canal WebSocket.
    \item Persistencia PostgreSQL sobre \textbf{13} tablas, con migraciones
    gestionadas v\'ia Alembic.
    \item Bus Redis Pub/Sub con dos bases de datos l\'ogicas (cach\'e + cola).
    \item Worker as\'incrono para c\'alculos pesados (RQ).
    \item Secuenciador para automatizaci\'on temporizada.
    \item Autenticaci\'on JWT con \texttt{bcrypt}.
    \item Sistema LIMS con plantas, experimentos y asignaciones N:M.
\end{itemize}

\subsection{Capa de Visualizaci\'on (React)}
\begin{itemize}[leftmargin=*]
    \item \textit{Single Page Application} con dashboard tiempo real.
    \item Control manual de actuadores.
    \item Planificador de secuencias multi-paso.
    \item Gesti\'on LIMS: plantas, experimentos, asignaciones,
    consulta de telemetr\'ia por planta.
\end{itemize}

\subsection{SDK Cient\'ifico (Python)}
\begin{itemize}[leftmargin=*]
    \item Cliente con cach\'e Parquet autom\'atica
    (\texttt{\textasciitilde/.demeter\_cache}).
    \item Pipeline ETL: \texttt{to\_dataframe} $\to$ \texttt{clean}
    $\to$ \texttt{fill\_gaps} $\to$ \texttt{resample}.
    \item C\'alculos agron\'omicos vectorizados:
    \begin{itemize}
        \item Termodin\'amica del aire h\'umedo: VPD, punto de roc\'io,
        bulbo h\'umedo (Stull, 2011), humedad absoluta, entalp\'ia (ASHRAE).
        \item Fenolog\'ia: GDD, \textit{chill hours},
        ET$_0$ (Hargreaves--Samani), DAP.
        \item Estad\'istica: Z-score, \'indice de Wallin (1962), media m\'ovil 24\,h.
    \end{itemize}
    \item Visualizaci\'on con Matplotlib/Seaborn (4 funciones gr\'aficas).
    \item Exportaci\'on CSV (UTF-8 BOM) y Excel multi-hoja.
    \item Suite completa en una sola llamada (\texttt{run\_full\_suite()}).
\end{itemize}

\section{Documentaci\'on y artefactos generados}

\begin{itemize}[leftmargin=*]
    \item Documentaci\'on t\'ecnica de arquitectura: \textbf{84~p\'aginas}
    en LaTeX, m\'as las 14.575 l\'ineas de documentaci\'on en Markdown del
    \'arbol \texttt{docs/}.
    \item \textbf{18} diagramas distintos de flujo, secuencia y
    arquitectura (compilados en variante clara y oscura).
    \item \textbf{9} cuadros de \textit{Decisi\'on de Dise\~no} con
    justificaci\'on y alternativas descartadas.
    \item Repositorio Git con \textbf{227} \textit{commits} y
    \textbf{27.025} l\'ineas de c\'odigo propio, excluidas dependencias y
    ficheros generados.
    \item Suite de \textbf{140 tests} automatizados repartidos en 20 ficheros:
    \begin{itemize}
        \item Firmware ESP32 (Unity/C++): 32 tests.
        \item Servicio Raspberry Pi (pytest): 37 tests.
        \item Backend FastAPI (pytest): 16 tests.
        \item Frontend (Vitest): 3 tests.
        \item SDK cient\'ifico (pytest): 52 tests.
    \end{itemize}
    \item Esquema de datos PostgreSQL versionado con Alembic.
    \item Modelos Pydantic compartidos entre todas las capas.
\end{itemize}

\section{Aprendizajes adquiridos}

Los miembros participantes han desarrollado competencias en:

\begin{itemize}[leftmargin=*]
    \item Dise\~no de protocolos de comunicaci\'on binarios.
    \item Programaci\'on embebida ESP32 + ESP-NOW.
    \item Programaci\'on as\'incrona en Python (\texttt{asyncio}).
    \item Backend moderno (FastAPI, asincron\'ia, Pydantic, SQLAlchemy as\'incrona).
    \item Frontend reactivo (React, WebSockets en cliente).
    \item DevOps b\'asico (Docker, despliegue cloud, gesti\'on de secretos).
    \item Documentaci\'on t\'ecnica profesional (LaTeX, diagramas).
    \item C\'alculos agron\'omicos vectorizados (NumPy/Pandas).
    \item Especificaci\'on y redacci\'on de decisiones de arquitectura
    (ADR) y de exploraciones acotadas (\textit{spikes}) como herramientas
    de trabajo, no como burocracia.
    \item Evaluaci\'on cr\'itica de la sostenibilidad de un proyecto propio
    y capacidad de proponer su cierre ordenado, que es en s\'i misma una
    competencia de ingenier\'ia.
\end{itemize}
